\chapter{\hybrid}
\label{ch:hybrid}
% Framing depends on the novelty decision (see Chapter 1 TODO). Whatever is
% decided, describe the model honestly: these components were ADDED and
% motivated by RFI morphology; Chapter 6 then tests whether they helped.

\section{Architecture}
% Code: src/hybrid_rfi_package/hybrid_model.py. Diagram needed (I can draw it).
% Same U shape as tf_unet. Differences, simply (you have the plain-language
% list from our chat):
%  - Residual blocks: each block learns a change on top of a shortcut.
%  - Multi-scale strip convolutions: 1xK and Kx1 kernels, K = 7, 11, 21,
%    look along long lines in time and frequency. Motivation: RFI is
%    often long thin streaks.
%  - ECA channel attention: re-weights feature channels.
%  - GroupNorm after convolutions.
%  - Raw logits output -- no ReLU on the output layer (avoids Ch. 4 trap).
%  - Same padding -- full 512x512 output.
% Size: base 8 -> 593,842 params; base 32 -> 9,304,186 [PART 6].
% tf_unet for comparison: 465,986.

\section{Training}
% - PyTorch. Loss = cross-entropy + Dice. Adam lr 1e-3, batch 1.
% - 700 steps x 40 epochs = 28,000 gradient steps (same budget as synthetic).
% - Best checkpoint and threshold chosen on 735 validation images.
% - Fixed-range normalisation from training images only.
% [source: experiments/lofar_hybrid.py, run metrics.json]

\section{Results on synthetic data}
% Keep short; synthetic is the easy case.
%  - 276x600: 0.9808 (base 32) [Where the trained models live table]
%  - 1024x265 v4: 0.9878 (base 8)
%  - Width sweep [PART 6]: base 4 0.9547 | base 8 0.9749 | base 16 0.9788 |
%    base 32 0.9812. 15.7x fewer parameters cost only 0.0063.
%    Base 4 genuinely breaks. (These are N=1 -- say so.)

\section{Results on LOFAR}
\subsection{Class weighting}
% - First run used class weights (56.9x on RFI): 0.6467 +/- 0.0079, but a
%   validation threshold of 0.9619 -- badly calibrated. [PART 13.6]
% - Without class weights: 0.6603 +/- 0.0040, threshold 0.2573.
%   Better in every seed, paired t = 6.03 on 2 dof. [PART 13.11]
% - Interesting: ROC fell (0.9801 -> 0.9378) while F1 rose. Ranking and
%   decision quality are different things. [PART 13.13]

\subsection{Final comparison}
% Table, max F1 (the paper's protocol) [PART 13.12]:
%  sigma-clip                  0.4103
%  tf_unet lr 1e-4 (ours)      0.5482 +/- 0.0139
%  AOFlagger                   0.5698
%  Mesarcik et al. U-Net       0.5876 +/- 0.0031
%  RFI-Net (published best)    0.5979
%  HybridRFINet base 8         0.6603 +/- 0.0040   (593,842 params)
% Also give pooled F1: 0.6561 +/- 0.0044 [PART 20 table].
% - +0.0624 over RFI-Net, t = 27.22 on 2 dof.
% - Beats its own teacher: trained on AOFlagger labels (0.5698), scores
%   0.66 against the human.
% - Base 32 at 9.3M params: 0.6571 (N=1, with class weights) -- no gain
%   from 15.7x more parameters [PART 13.7].

\section{Checking the result}
% PART 14 -- an adversarial self-audit. This makes the result credible.
%  - Metrics recomputed by independent code: differences <= 3.2e-5.
%  - Leakage: 0 train/val images in the test set. Normalisation range
%    computed from training images only.
%  - Honest thresholding costs 0.0064 vs choosing on test.
%  - Train/val/test max F1 0.7112 / 0.6820 / 0.6592 -- mild overfitting.
%  - Not copying AOFlagger: agreement with it only 0.7430; finds 16,485
%    true-RFI pixels AOFlagger missed, misses 7,345 it found.
%  - Not a trivial "flag the bad columns" rule: that scores 0.0645.
%  - Bootstrap over the 109 test images (POOLED F1, 2000 resamples):
%    3-seed mean 0.6561, 95% CI [0.6154, 0.6923]. RFI-Net's 0.5979 lies
%    outside; 99.5% of resamples beat it. Note the CI is +/-0.04 wide --
%    the 109-image test set is the biggest source of uncertainty.
